\uuid{YUup}
\exo7id{4725}
\titre{exo7 4725}
\auteur{quercia}
\organisation{exo7}
\datecreate{2010-03-16}
\isIndication{false}
\isCorrection{false}
\chapitre{Topologie}
\sousChapitre{Topologie de la droite réelle}
\cpge{mpsi-mp2i.nombres-reels-suites.d, pcsi.nombres-reels-suites.c, ptsi.nombres-reels-suites.c, mpsi-mp2i.nombres-reels-suites.g, bcpst1.suites-reelles, mp-mpi.topologie-espaces-normes.d, pc.espaces-vectoriels-normes.d, psi.espaces-vectoriels-normes.d, pt.fonctions-deux-trois-variables.a, mp-mpi.topologie-espaces-normes.e, pc.espaces-vectoriels-normes.e, psi.espaces-vectoriels-normes.e, pt.fonctions-deux-trois-variables.b}
\module{Analyse}
\niveau{L2}
\difficulte{}

\contenu{
\question{
Soit $(u_n)$ une suite r{\'e}elle telle que $u_{n+1} - u_n \xrightarrow[n\to\infty]{} 0$
et $u_n \xrightarrow[n\to\infty]{}+\infty$.
D{\'e}montrer que la suite $(\exp(iu_n))$ est dense dans $\mathbb{U}$.
}
}
